%! Characteristics of Quadratic Functions — Unit 2, Lesson 2.0 — Slides (Beamer)
% ELEVEN FRAMES, in the profile's order (LESSON_SHAPE.md §1): title → learning targets (+ how
% today runs, 5/35/10/10) → warm-up (ending in "hold on to this") → hook (the vote, LEFT
% UNRESOLVED) → four notes frames, two per section (rows 1–2 = section 1, rows 3–4 = section 2;
% the crux frame flagged \sectionlabel[redacc]{}) → Guided Practice → debrief → close & start the
% homework. There is no group activity and no individual-practice launch frame: the homework IS
% the individual practice, started in class. \CourseName does not exist in beamer — the course
% name is written literally. Header macro is \forestheader{Title}.
\documentclass[aspectratio=169, 11pt]{beamer}
\usepackage{algebra2-beamer}   % provides \forestheader{} and \sectionlabel[color]{}
% Coordinate grid: {xmin}{xmax}{ymin}{ymax}
\newcommand{\lgrid}[4]{%
  \draw[step=1, linegray!40, very thin] (#1,#3) grid (#2,#4);
  \draw[->, thick, charcoal] (#1,0)--(#2,0) node[below right, font=\tiny]{$x$};
  \draw[->, thick, charcoal] (0,#3)--(0,#4) node[left, font=\tiny]{$y$};}
% Projectile grid: ONE GRID UNIT UP = 16 FEET, so the plotted curve is h/16 = -t^2 + 2t + 3.
\newcommand{\pgrid}{%
  \draw[step=0.5, linegray!40, very thin] (0,0) grid (3.5,4.5);
  \draw[->, thick, charcoal] (0,0)--(3.9,0) node[below right, font=\tiny]{$t$ (s)};
  \draw[->, thick, charcoal] (0,0)--(0,5.0) node[left, font=\tiny]{$h$ (ft)};
  \foreach \x in {1,2,3}{\node[below, font=\tiny, charcoal] at (\x,0){$\x$};}
  \foreach \y/\lab in {1/16,2/32,3/48,4/64}{\node[left, font=\tiny, charcoal] at (0,\y){$\lab$};}}
\newcommand{\fdot}[3]{\fill[#1] (#2,#3) circle (3.0pt);}

\begin{document}

% 1. TITLE SLIDE (CourseName is not defined in beamer, so the name is literal)
{
\setbeamercolor{background canvas}{bg=forest}
\begin{frame}[plain]
  \vspace{2.8cm}\hspace{0.55cm}%
  \begin{minipage}{0.9\textwidth}
    {\color{goldacc}\bfseries\small LESSON 2.0}\\[0.5cm]
    {\color{white}\bfseries\LARGE Characteristics of Quadratic Functions}\\[1.2cm]
    {\color{forestmid}\small Algebra 2~$\cdot$~Unit 2: Quadratic Functions}
  \end{minipage}
\end{frame}
}

% 2. LEARNING TARGETS
\begin{frame}[t, plain]
  \forestheader{Today's targets}
  \sectionlabel{What you will be able to do}
  \begin{itemize}\setlength{\itemsep}{5pt}
    \item Find the \textbf{vertex} of a \textbf{parabola} and say whether it is a
          \textbf{maximum} or a \textbf{minimum} point.
    \item State that \textbf{maximum or minimum value} --- an \emph{output}, with units --- and
          separately say \emph{where} it occurs.
    \item Write the \textbf{axis of symmetry} $x=h$, and use symmetry to pair equal outputs and
          to find the vertex from a table or from the two zeros.
    \item Read the \textbf{zeros}, the \textbf{$y$-intercept}, the \textbf{increasing} and
          \textbf{decreasing} intervals, the \textbf{domain} and \textbf{range}.
    \item Describe \textbf{end behavior}: $a>0$ both ends rise, $a<0$ both ends fall.
  \end{itemize}
  \vspace{4pt}
  \begin{block}{How today runs}
    Warm-up (5) $\to$ guided notes and one graph read together (35) $\to$ debrief (10) $\to$
    \textbf{start the homework, on your own} (10). Everything today is \emph{reading} a graph ---
    no solving, no factoring.
  \end{block}
\end{frame}

% 3. WARM-UP
\begin{frame}[t, plain]
  \forestheader{Warm-Up: three things you already know}
  \sectionlabel{5 minutes --- alone, then check with a neighbour}
  \begin{enumerate}\setlength{\itemsep}{7pt}
    \item \textbf{Evaluate the rule.} $f(x)=x^2-2x-3$ at $-1$, $0$, $1$, $3$. Watch the
          negatives. \; (You get $0,\,-3,\,-4,\,0$ --- hold on to those.)
    \item \textbf{Read the line.} For $y=-2x+4$: the $x$-intercept, the $y$-intercept,
          increasing or decreasing, and the domain.
    \item \textbf{A graph that turns around.} Increasing until? Decreasing after? Highest point?
  \end{enumerate}
  \vspace{6pt}
  \begin{block}{Hold on to this}
    The highest point of that graph is $(2,4)$. \textbf{Is the ``highest'' answer the $x$-value,
    the $y$-value, or both?} Keep that question --- it is the whole lesson.
  \end{block}
\end{frame}

% 4. HOOK — left unresolved
\begin{frame}[t, plain]
  \forestheader{Hook: a ball off a 48-foot ledge}
  \sectionlabel[goldacc]{$h(t)=-16t^2+32t+48$ \quad ($t$ in seconds, $h$ in feet)}
  \begin{columns}[T]
    \begin{column}{0.46\textwidth}
      \centering
      \begin{tikzpicture}[scale=0.62]
        \pgrid
        \draw[very thick, forest, domain=0:3, samples=80, smooth] plot (\x, {-\x*\x+2*\x+3});
        \fdot{forest}{0}{3}
        \fdot{forest}{3}{0}
        \fdot{redacc}{1}{4}
      \end{tikzpicture}
    \end{column}
    \begin{column}{0.50\textwidth}
      \begin{itemize}\setlength{\itemsep}{4pt}
        \item It starts at $h(0)=48$ ft --- the top of the ledge.
        \item It turns \emph{once}, at the point $(1,64)$.
        \item It hits the ground at $t=3$ s.
        \item \textbf{One grid unit up is $16$ feet} --- on a plain unit grid this curve would be
              a page-tall spike.
      \end{itemize}
    \end{column}
  \end{columns}
  \vspace{2pt}
  \begin{block}{Vote --- and we do not settle it yet}
    \textbf{The ball's maximum is: circle one --- $1$ or $64$.} \quad And \emph{what are the units
    of your answer?} Hands up for each. Keep your vote; we come back to it.
  \end{block}
\end{frame}

% 5. NOTES ROW 1 — THE TURNING POINT (section 1a)
\begin{frame}[t, plain]
  \forestheader{Notes: a parabola turns exactly once}
  \sectionlabel{notes, The turning Point $\cdot$ I do problem 1, you do problem 2}
  \begin{itemize}\setlength{\itemsep}{5pt}
    \item A \textbf{parabola} is the graph of $f(x)=ax^2+bx+c$, $a\ne0$. The \emph{sign of $a$}
          says which way it opens: $a>0$ up, $a<0$ down.
    \item The \textbf{vertex} is the point where it turns: the \textbf{minimum} point when it
          opens up, the \textbf{maximum} point when it opens down.
    \item Our ball: $a=-16<0$, so it opens down --- and the turning point $(1,64)$ is the
          \emph{highest} point on the curve.
  \end{itemize}
  \vspace{3pt}
  \begin{block}{One point does most of the work}
    Because it turns only once, the vertex plus the direction it opens already tell you the
    maximum or minimum, the mirror line, the rising and falling parts, the range, and both ends.
  \end{block}
\end{frame}

% 6. NOTES ROW 2 — THE MIRROR LINE (section 1b)
\begin{frame}[t, plain]
  \forestheader{Notes: the mirror line}
  \sectionlabel{notes, The mirror Line $\cdot$ the unit anchor curve (problems 3--4)}
  \begin{columns}[T]
    \begin{column}{0.44\textwidth}
      \centering
      \begin{tikzpicture}[scale=0.34]
        \lgrid{-2.5}{4.5}{-5}{5.5}
        \draw[dashed, very thick, redacc] (1,-5)--(1,5.2);
        \draw[very thick, forest, domain=-1.9:3.9, samples=80, smooth] plot (\x, {\x*\x-2*\x-3});
        \fdot{forest}{1}{-4}
      \end{tikzpicture} \\[2pt]
      $f(x)=x^2-2x-3=(x-1)^2-4$
    \end{column}
    \begin{column}{0.52\textwidth}
      \begin{itemize}\setlength{\itemsep}{4pt}
        \item \textbf{Axis of symmetry}: the vertical line $x=h$ through the vertex --- here
              $x=1$.
        \item Inputs the same distance either side give \textbf{equal outputs}:
              $f(0)=f(2)=-3$, \; $f(-1)=f(3)=0$.
        \item The parent $y=x^2$ is the plain case: $f(-2)=f(2)=4$ --- \emph{even symmetry}.
        \item So the axis is midway between the zeros: $\frac{-1+3}{2}=1$.
      \end{itemize}
    \end{column}
  \end{columns}
  \vspace{3pt}
  \begin{block}{Remember this curve}
    $f(x)=x^2-2x-3$ is the Unit 2 anchor --- vertex $(1,-4)$, axis $x=1$, zeros $-1$ and $3$,
    $y$-intercept $(0,-3)$. It comes back in 2.1 through 2.5 wearing different costumes.
  \end{block}
\end{frame}

% 7. NOTES ROW 3 — READING THE WHOLE GRAPH (section 2a)
\begin{frame}[t, plain]
  \forestheader{Notes: reading the whole graph}
  \sectionlabel{notes, Reading the Whole Graph $\cdot$ problems 5--7, the feature card}
  \begin{columns}[T]
    \begin{column}{0.50\textwidth}
      \begin{itemize}\setlength{\itemsep}{4pt}
        \item \textbf{Zeros} ($x$-intercepts): $-1$ and $3$. A parabola has \emph{two, one, or
              none} --- never three.
        \item \textbf{$y$-intercept}: the point $(0,-3)$.
        \item \textbf{Decreasing} on $(-\infty,1]$, \textbf{increasing} on $[1,\infty)$ --- the
              axis is where it turns, so the axis is the split.
      \end{itemize}
    \end{column}
    \begin{column}{0.46\textwidth}
      \begin{itemize}\setlength{\itemsep}{4pt}
        \item \textbf{Domain}: all real numbers --- always, for every quadratic.
        \item \textbf{Range}: $y\ge k$ if it opens up, $y\le k$ if it opens down, where $k$ is the
              vertex's output. Here $y\ge-4$.
      \end{itemize}
    \end{column}
  \end{columns}
  \vspace{3pt}
  \begin{block}{Two answers that look alike but are not}
    A \textbf{zero} is an $x$-value: $-1$. A \textbf{$y$-intercept} is a \emph{point}: $(0,-3)$.
    Say which kind of thing each answer is before you write it down.
  \end{block}
\end{frame}

% 8. NOTES ROW 4 — HOW HIGH vs. WHERE — THE CRUX (section 2b)
\begin{frame}[t, plain]
  \forestheader{Notes: how high, versus where}
  \sectionlabel[redacc]{notes, How high vs.\ Where $\cdot$ the crux, problem 8 --- the value is an OUTPUT}
  \begin{itemize}\setlength{\itemsep}{4pt}
    \item ``The ball's maximum is $1$''? \textbf{No.} $1$ answers \emph{when} --- it is the
          \emph{input}, measured in seconds.
    \item The maximum \textbf{value} is $64$, measured in \emph{feet}: an \emph{output}. The
          vertex $(1,64)$ answers both questions at once, in that order --- \emph{where}, then
          \emph{how high}.
    \item \textbf{End behavior}: $a>0$, both ends rise; $a<0$, both ends fall. A parabola's two
          ends always go the \emph{same} way --- a line's go opposite ways.
  \end{itemize}
  \vspace{3pt}
  \begin{block}{Where the two answers disagree --- and they almost always do}
    They agree only by accident. Here $1\ne64$: one is a time, the other a height, and no amount
    of arithmetic turns seconds into feet. So the question decides the answer: \emph{``what is the
    maximum?''} wants $64$ ft; \emph{``when does it happen?''} wants $t=1$ s. \textbf{Always name
    the units.}
  \end{block}
\end{frame}

% 9. GUIDED PRACTICE — WE DO
\begin{frame}[t, plain]
  \forestheader{Guided Practice: we read this one together}
  \sectionlabel[goldacc]{You hold the pen --- problems 10--13 on $g(x)=-(x+2)^2+9$}
  \begin{columns}[T]
    \begin{column}{0.36\textwidth}
      \centering
      \begin{tikzpicture}[scale=0.26]
        \lgrid{-6.5}{2.5}{-4}{10.5}
        \draw[dashed, semithick, redacc] (-2,-4)--(-2,10.2);
        \draw[very thick, forest, domain=-5.6:1.6, samples=80, smooth]
          plot (\x,{-(\x+2)*(\x+2)+9});
        \fdot{forest}{-2}{9} \fdot{redacc}{-5}{0} \fdot{redacc}{1}{0} \fdot{navy}{0}{5}
      \end{tikzpicture}
    \end{column}
    \begin{column}{0.60\textwidth}
      \begin{enumerate}\setlength{\itemsep}{3pt}
        \item[\textbf{10.}] Which way does it open? Vertex as a \emph{point}, axis as an
              \emph{equation}.
        \item[\textbf{11.}] Zeros and $y$-intercept --- then check both zeros against the axis.
        \item[\textbf{12.}] Increasing where, decreasing where, domain, range, end behavior.
        \item[\textbf{13.}] The \emph{maximum value}, and \emph{where} it occurs. A classmate
              answers ``$-2$'': what have they confused?
      \end{enumerate}
    \end{column}
  \end{columns}
  \vspace{2pt}
  \begin{block}{The questions I will ask}
    ``How far is $-5$ from the axis? How far is $1$?'' \quad ``Where does it turn --- so where
    does increasing stop?'' \quad ``Is $9$ an input or an output?''
  \end{block}
\end{frame}

% 10. DEBRIEF
\begin{frame}[t, plain]
  \forestheader{Debrief: pens down, papers up}
  \sectionlabel[redacc]{10 minutes --- aloud, nothing to fill in}
  \begin{enumerate}\setlength{\itemsep}{5pt}
    \item \textbf{The vertex $(1,64)$, out loud, twice.} The maximum value is $64$ \emph{feet};
          it happens at $t=1$ \emph{second}. That settles the vote --- by naming the units.
    \item \textbf{Problem 13.} The maximum value of $g$ is $9$; $-2$ is \emph{where}. Whoever
          wrote $-2$ gave an input where an output was asked for --- same mistake, new curve.
    \item \textbf{Symmetry did the work.} $-5$ and $1$ are each $3$ from $x=-2$; averaging the
          zeros finds the axis; $f(0)=f(2)=-3$ on the anchor. One idea, three uses.
    \item \textbf{Domain and end behavior.} The \emph{function} runs over all real numbers and
          both ends go the same way; only the \emph{ball} stops when it lands.
  \end{enumerate}
  \vspace{3pt}
  \begin{block}{Say it without the notes}
    ``The vertex $(1,64)$ means the greatest height is $64$ feet, and it happens one second after
    launch.''
  \end{block}
\end{frame}

% 11. CLOSE & START HOMEWORK
\begin{frame}[t, plain]
  \forestheader{Close --- and start the homework, on your own}
  \sectionlabel{10 minutes $\cdot$ silent $\cdot$ the Homework page in your packet}
  You walked in able to read a \emph{line}. You walk out able to read a \emph{parabola} --- where
  it turns, how high that is, where it crosses, where it rises and falls, and what both ends do
  --- and to say which of those answers is an input and which is an output.
  \vspace{5pt}
  \begin{block}{Homework --- scored, due the first class after two study halls}
    We read item~1's vertex together right now; then you keep going, alone. Six items: one graph
    read whole, a contrast pair, a table and its second differences, the special cases, a punt,
    and one multiple-choice item. Whatever is left is due the first class after two study halls
    --- I will announce the date in class and on TurtleNet.
  \end{block}
  \vspace{4pt}
  \textbf{Next:} Lesson 2.1 --- the same curve in three different costumes, and the graph built
  \emph{from} the equation instead of read off a picture.
\end{frame}

\end{document}
